\section{Uniform constants and an explicit finite-record inverse}
\label{sec:v78-uniform}

The finite-cover theorem gives a protocol for each table.  We now
separate its finite design parameters from the estimates of the inverse.
This makes the long-prefix cost visible and permits a constructive
finite-record output.  It does not assert a bound for the number of
patches on every positive-curvature table from curvature alone.

\subsection{Primitive geometric bounds}
Suppose every flight in the finite protocol has length in
$[\ell_-,\ell_+]$, every incidence cosine is at least $\nu_->0$,
every scalar chart speed belongs to $[v_-,v_+]$, and curvature belongs
to $[\kappa_-,\kappa_+]$, with positive lower endpoints.  Let $N_* $
be the largest number of flights in any measured branch.  Set
\begin{align*}
 b_*&=v_+^2/\ell_-,&
 b_0&=v_-^2\nu_-^2/\ell_+,&
 h_0&=2\kappa_-\nu_-v_-^2,\\
 D&=2v_+^2/\ell_-+2\kappa_+v_+^2,&
 u_*&=b_0(b_0/D)^{N_*-1},&
 U_*&=\max(b_*,b_*^2/h_0).
\end{align*}
All intervals and derivatives are taken on fixed closed gates with
slightly larger regular domains.

\begin{proposition}[Explicit twist and acceptance margins]
\label{prop:v78-geometric-bounds}
On this class every measured action satisfies
\begin{equation}\label{eq:v78-billiard-twist}
                    u_*\le |W_{st}|\le U_*.
\end{equation}
At a prefix match the positive Schur complement is at most $D$, and
the magnitude of the matching derivative is at least $u_*^2/D$.
If the normalized preparation density is at least $\rho_->0$, retention
is at least $e_->0$, the smallest clock residual is $r_->0$, and the
smallest gate area is $q_->0$, then the retained-success probability
of every attempted setting is at least
\begin{equation}\label{eq:v78-derived-acceptance}
                         p_*=e_-\rho_-u_*r_-q_->0.
\end{equation}
\end{proposition}
\begin{proof}
An edge coupling has absolute value between $b_0$ and $b_*$ by
\eqref{eq:v77-edge-twist}.  The interior Hessian is bounded below by
$h_0 I$ by its second-variation formula.  Its diagonal entries are at
most $D$: the two projected kinetic terms contribute at most
$2v_+^2/\ell_-$ and the curvature term at most
$2\kappa_+v_+^2$.  Parameter acceleration along the tangent cancels
at an interior stationary reflection.

For a branch of $N$ flights the corner-cofactor identity gives
$|W_{st}|=\prod_{i=0}^{N-1}|b_i|/\det H_{N-1}$.
Hadamard's inequality bounds the determinant by $D^{N-1}$.
Since $b_0/D\le1$, the stated lower bound follows for $N\le N_*$.
For $N\ge2$, the endpoint Schur formula and
$\|H_{N-1}^{-1}\|\le h_0^{-1}$ give the upper bound
$b_*^2/h_0$; for $N=1$ use $b_*$.  The last-site Schur complement
of the extended Hessian is positive and no larger than its original
diagonal entry, hence at most $D$.  Apply~\eqref{eq:v78-root-sign}.
Finally the actual retained mass is the integral of
$e\rho|W_{st}|(T-W)$ over the gate, by the physical flow-box calculation.
Integrating the indicated lower bounds proves~\eqref{eq:v78-derived-acceptance}.
Here $\rho$ is a normalized preparation density; an arbitrary unnormalized
weight is not silently turned into a success probability.
\end{proof}

\begin{proposition}[An anchor from three corner tests]
\label{prop:v78-anchor-bound}
Suppose $|W_{st}|\ge u>0$ on a rectangle of side lengths $d_s,d_t$,
$a_-\le a\le a_+$ there, and $0<r_-\le T_j-W\le r_+$.
Take $z_0$ to be one corner.  Among the other three corners one has the
three-clock anchor denominator bounded below by
\begin{equation}\label{eq:v78-anchor-bound}
 d_*:=\frac{a_-}{a_+}\frac{u d_s d_t}{3}
 \frac{(T_2-T_1)(T_3-T_1)(T_3-T_2)}{r_+^3}>0.
\end{equation}
The corner with largest observed denominator is therefore usable.  The
rectangle can lie strictly inside the original normalization gate.
\end{proposition}
\begin{proof}
The mixed derivative has constant sign on the rectangle.  Its integral
is the four-corner difference of $W$ and has magnitude at least
$u d_s d_t$.  Writing this difference as a signed sum of the other
three corner values minus $W(z_0)$ shows that one has
$|W(z)-W(z_0)|\ge u d_s d_t/3$.
The identity in the proof of Theorem~\ref{thm:v77-clock} gives
\[
 \left|\sum_jc_jg_j(z)\right|
 =\frac{a(z)}{a(z_0)}|W(z)-W(z_0)|
   \frac{(T_2-T_1)(T_3-T_1)(T_3-T_2)}
        {\prod_j(T_j-W(z_0))}.
\]
The bounds prove the claim.  Restricting functions without changing
their original normalizations leaves all ratios unchanged.
\end{proof}

\subsection{One finite design on a bounded neighborhood}
Fix a reference table and one finite atlas constructed in
Theorem~\ref{thm:v77-finite-atlas}.  The acquisition list consists only
of its branch words, scalar gates, and clocks.  Clock anchors, root
subbrackets, distance anchors and registration witnesses may be used
in the proof of validity of the class, but are not provided afresh for
an unknown member.  The procedure uses the recovered actions and
strict tests to find these witnesses.

\begin{theorem}[Reference-uniform finite protocol]
\label{thm:v78-uniform}\label{thm:v77-stability}
For every such reference atlas there is a positive finite-smooth
neighborhood on which its one fixed acquisition list reconstructs all
tables.  After imposing bounds for the finitely required derivatives of
geometry and source and positive source and retention floors, the same
list has fixed $M,N_*,u_*,p_*,d_*>0$, fixed distance and overlap
singular-value margins, and a uniform inverse error bound
\begin{equation}\label{eq:v78-geometric-stability}
 d_k(\mathcal T,\widetilde{\mathcal T})
 \le C_k\max_{\gamma,j}
       \|f_{\gamma,j}-\widetilde f_{\gamma,j}\|_{C^{k+2}}
\end{equation}
for sufficiently close realizable law tuples.  Here $d_k$ is the
maximum boundary $C^k$ error and lattice error modulo one common
Euclidean motion.  This statement is uniform on the prescribed
neighborhood, not on the union of all such neighborhoods.
\end{theorem}
\begin{proof}
The strict inequalities in a reference finite atlas are finitely many
compact-domain clearance, incidence, residual, root and rank conditions.
Finite-branch implicit dependence and continuity preserve them in a
finite-smooth neighborhood.  If $\eta_i>0$ are their reference margins
and $L_i$ bound their variations per unit change of the embeddings on a
larger bounded neighborhood, then any radius
\[
                  0<\delta<\min_i\eta_i/[2(1+L_i)]
\]
preserves half of every margin.  Such finite $L_i$ follow by implicit
differentiation using the interior Hessian inverse, followed by the
finite matrix and evaluation operations defining the certificates.
They can be bounded recursively from finite derivative bounds and the
Hessian lower bound.  The reference certificate, rather than a universal
numerical atlas bound, specifies $M$ and $N_*$.  The explicit triangular
support-function family in Proposition~\ref{prop:v77-example} supplies
nonempty neighborhoods of this kind with nonanalytic smooth members.

The preceding propositions now give uniform twist, success and clock
anchor bounds.  The fixed distance and overlap matrices retain their
singular-value margins.  Recover actions without differentiation; recover
roots and distances with two derivatives of loss; recover boundary
patches by the finite metric solve and trilateration; and register along
a fixed finite spanning tree.  Each operation has a uniform local bound
on these sets.  The root operation may be implemented on its uniquely
matched domain without a supplied choice among roots.  Composing the
bounds gives~\eqref{eq:v78-geometric-stability}, including the lattice
columns recovered by differences of translated intervals.  This is a
composition of explicit inverses, not an appeal to continuity of an
unspecified compact inverse.
\end{proof}

In particular the matching inverse retains the possible exponential
factor $D/u_*^2$ in $N_*$.  Formula~\eqref{eq:v78-derived-acceptance}
charges the corresponding possible rarity of successes.  Distance rank
margins and registration depth are also retained.  The theorem neither
removes grazing or near-degenerate finite certificates nor asserts
nonidentifiability of conic arcs.

\subsection{From estimated laws to an actual table}
For a whole-table protocol use the fixed scalar-circle label convention.
Choose a fixed smooth partition of unity on each circle, subordinate
to slightly smaller intervals of the finite cover.  Each registered
reconstructed patch is a smooth function on the larger interval.
Blend the patches with this partition to obtain a closed curve for each
quotient obstacle.  Use the recovered two lattice vectors to generate
all its translates.  All registration operations use polar
orthogonalization of the invertible maps determined by overlap triples,
so are defined smoothly near the exact data even when noisy triples
are not exactly congruent.

\begin{theorem}[Constructive finite-record reconstruction]
\label{thm:v78-records}
Fix a bounded neighborhood as in Theorem~\ref{thm:v78-uniform}, an
integer $k\ge2$, and $s\ge1$.  Suppose the conditional densities
have bounded $C^{k+2+s}$ norms on larger gates.  Make $B$ independent
attempts at each of the $M$ settings and charge every attempt.  Retained
coordinates may have arbitrary errors of magnitude at most $\delta$.
There is a reconstruction using kernel smoothing, explicit clock
formulas, scalar root solving, finite linear algebra and patch blending
such that, for sufficiently large $Bp_*$ and small recording error,
\begin{equation}\label{eq:v78-record-rate}
 d_k(\widehat{\mathcal T},\mathcal T)
 \le C\left\{
 \left(\frac{\log(C_0MB/\zeta)}{Bp_*}\right)^{s/[2(k+2+s)+2]}
       +\delta^{s/(k+s+5)}\right\}
\end{equation}
with probability at least $1-2\zeta$.  On this event the output is a
smooth strictly convex disjoint periodic table, not merely a formal
list of inconsistent distance kernels.  The total preparation budget
is $MB$.  For $k=0,1$ the same constructive conclusion follows from
the $k=2$ bound; no sharper low-order rate is inferred from this argument.
\end{theorem}
\begin{proof}
Set $m=k+2$ and $n=\lfloor Bp_*/2\rfloor$.  A Chernoff bound and
union bound show that every setting has at least $n$ retained successes
except on an event of probability at most $M e^{-Bp_*/8}$.  The first
$n$ successes in each independent preparation sequence are coupled to
$n$ independent draws from the conditional law.  This couples the
capped experiment without assuming independence conditional on meeting
its cap.  For $Bp_*\ge C\log(M/\zeta)$, cap failure is at most
$\zeta$.

Use a smooth compactly supported order-$s$ kernel on each larger gate.
For each derivative of order at most $m$, Taylor expansion gives bias
$Ch^s$; the summand has envelope $Ch^{-m-2}$, variance
$Ch^{-2m-2}$, and Lipschitz bound $Ch^{-m-3}$.  Bernstein's inequality
on a grid of mesh $n^{-2}h^{m+3}$, followed by a union bound over
settings and derivatives, therefore gives, for $h\ge n^{-1}$,
\[
 e_n\le C\left(h^s+
 \sqrt{\frac{L_n}{nh^{2m+2}}}+\frac{L_n}{nh^{m+2}}
       +n^{-2}+\frac{\delta}{h^{m+3}}\right),\qquad
 L_n=\log(C_1Mn/\zeta),
\]
with probability $1-\zeta$.  The last term is the deterministic mean
value bound for moving every kernel center by at most $\delta$; the
recording errors need not be independent.  The derivative bound is on
inner gates, so no smooth zero extension at the boundary is assumed.
This is the elementary argument underlying the retained density lemma,
not a new optimal kernel rate; see also~\cite{GineGuillou2002}.
Choosing
\[
 h\asymp\max\{(L_n/n)^{1/[2(m+s)+2]},\delta^{1/(m+s+3)}\}
\]
makes all displayed terms bounded by the radius in
\eqref{eq:v78-record-rate}.

For this fixed-class procedure, retain a finite list of rational
certificate candidates with strict margins on the reference neighborhood.
Use only candidates passing half of those margins.  This list is fixed
once in the reconstruction rule; it is not a new mark supplied for an
unknown member.  It prevents an unrestricted certificate search from
selecting an arbitrarily ill-conditioned matrix.

Apply the clock inverse on a usable corner chart.  Small error preserves
its denominator; choosing among the three tested corners does not require
an unknown true anchor.  On compact matched patches the true root is
simple with a uniform margin.  Small $C^2$ action error preserves the
endpoint signs and a strictly decreasing neighborhood of that root.
On the compact complement of a small root neighborhood the exact matching function is bounded away from zero; after shrinking the reference class this separation is uniform.  Small perturbations therefore introduce no additional roots there.  Bisection and implicit differentiation give the smooth continuation.
The finite distance solves, positive matrix square roots, trilateration
and overlap registration likewise extend to sufficiently small
nonmodel perturbations.  The resulting registered patches are within
$Ce_n$ in $C^k$ of the true patches in one common gauge.  Blending is a
bounded linear operation in these fixed charts, and for exact patches
it returns the true closed boundary.  Thus its output has the same
$C^k$ error.  Lattice recovery is a bounded difference of registered
point evaluations.

For completeness, actual geometric validity is open in this norm.
A compact smooth embedded reference curve has a tubular neighborhood;
sufficiently small $C^1$ perturbations remain embedded and regular,
and small $C^2$ perturbations preserve positive curvature.  The lattice
remains invertible.  Disjointness of the finitely many nearby translated
obstacles follows from their positive separation.  For all remaining
translates, a common diameter bound and a lower lattice singular value
exclude intersection.  These margins can be made uniform on a smaller
reference neighborhood.  Hence the blended output is an actual periodic
table when $e_n$ is small.  In the explicit finite-horizon family the
contained-disk margin also preserves finite horizon.  Off the stated
event the procedure may return a fixed reference table whenever a
required inverse is undefined; no off-event accuracy is claimed.

The construction never enumerates candidate smooth tables or minimizes
distance to an unspecified observation image.  Root solving may be done
to any certified tolerance, adding its propagated numerical tolerance
to the displayed error.  No polynomial arithmetic-time bound or minimax
claim is made.  Combining the cap and estimation events proves the
probability assertion.
\end{proof}
